\chapter{Arbitrage de convertibles}\label{ch:s2:convertible-arbitrage}

Acheter la convertible, vendre l'action à découvert, et la position est longue de volatilité et de crédit, et courte de l'emprunt. Au premier trimestre 2005, les investisseurs
ont retiré plus de 20~\% du capital des fonds d'arbitrage de convertibles. Les fonds ont vendu des convertibles pour les rembourser, et les cours sont tombés sous les
valeurs théoriques, de 2.7~\% en médiane à la mi-mai. Les fonds ont perdu 7.2~\% de janvier à mai, soit à peu près ce que coûte un creusement de la décote de 2.7~\% à l'effet de levier typique de trois
pour un. En 2008, les prime brokers qui prêtaient aux hedge funds ont réduit l'effet de levier qu'ils accordaient. Sur le marché synthétique de ce chapitre, un portefeuille couvert en delta de
trente convertibles, à effet de levier de trois pour un, rapporte 9.8~\% par an sur son capital en temps ordinaire. Il perd 11.6~\% dans un épisode planté de
vente forcée de trois mois, dont 12.6 points viennent du creusement de la décote et 6.9 du crédit, puis il récupère 22.0~\% dans les six mois qui suivent. La construction est
\texttt{firm.convarb}.

\section{Ce que détient le porteur d'une convertible}

Une obligation convertible (livre~5, chapitre~21) est une obligation d'entreprise que le porteur peut échanger contre un nombre fixe d'actions. C'est une obligation simple, plus une option
sur les actions, moins ce que retire le rachat anticipé de l'émetteur. L'ensemble est exposé au défaut de l'émetteur. Le modèle de valorisation du livre~5 l'évalue sur une grille selon le cours de l'action
avec un taux de défaut instantané qui augmente lorsque l'action baisse~: un lien action-crédit.

L'obligation du chapitre a cinq ans de maturité, un coupon de 2~\%, et se convertit en 2.5 actions (un prix de conversion de 40). Elle est rachetable au bout de deux ans si
l'action vaut 30~\% de plus que le prix de conversion. Le taux de défaut est de 3~\% par an pour un cours de l'action de 40, et il augmente proportionnellement à $40/S$, avec un taux de recouvrement de 40~\%. À
l'émission, l'action valant 40, le modèle la valorise à 117.7. Sa valeur de conversion est de 100, et son delta est de 1.80 action par obligation
(la \cref{fig:s2:convertible-arbitrage:profile}). Pour un cours de l'action de 80, elle se négocie près de sa valeur de conversion (204.8 contre 200) avec un delta de 2.40.

\begin{definition}[Convertible «~busted~»]\label{def:s2:convertible-arbitrage:busted}
Une \emph{convertible «~busted~»}\index{convertible « busted »} est une convertible dont le cours de l'action est tombé très en dessous de son prix de conversion, de sorte que
l'option de conversion ne vaut presque rien et que l'obligation se négocie sur son crédit~; sa sensibilité résiduelle à l'action provient surtout du lien entre le cours
de l'action et le risque de défaut de l'émetteur.
\end{definition}

Pour un cours de l'action de 15, l'obligation est «~busted~». Elle vaut 81.2, en dessous des 87.34 de l'obligation simple au spread de crédit de l'émission, parce que le taux de défaut a augmenté
avec la baisse. Son delta reste de 1.27 action, une élasticité à l'action de seulement 0.23, et presque tout ce delta est du crédit.

\begin{omfigure}
\begin{tikzpicture}
\begin{axis}[width=10cm,height=6cm,xlabel={\small cours de l'action},ylabel={\small valeur pour 100 de nominal},tick label style={font=\footnotesize},
  xmin=8,xmax=100,ymin=0,ymax=260,grid=major,grid style={gray!25},table/col sep=comma,
  legend style={font=\scriptsize,at={(0.5,-0.25)},anchor=north},legend columns=3]
\addplot[omThm,thick] table[x=share,y=value]{figdata/strategies-2/08-convertible-arbitrage/profile.csv};
\addplot[omDef,thick,dashed] table[x=share,y=conversion]{figdata/strategies-2/08-convertible-arbitrage/profile.csv};
\addplot[gray,thick,dotted] table[x=share,y=floor]{figdata/strategies-2/08-convertible-arbitrage/profile.csv};
\legend{convertible (modèle),valeur de conversion,obligation simple au spread de l'émission}
\end{axis}
\end{tikzpicture}

\omcaption{La convertible de cinq ans du chapitre à l'émission, selon le cours de l'action~: valeur calculée par le modèle de valorisation du livre~5 avec un taux de défaut action-crédit, valeur de conversion
(2.5 actions), et l'obligation simple actualisée au spread de crédit de l'émission. Pour un cours de l'action de 19 et en dessous, l'obligation vaut moins que l'obligation simple
car le taux de défaut augmente avec la baisse. Données~: \texttt{s2\_convarb.profile}.}\label{fig:s2:convertible-arbitrage:profile}
\end{omfigure}

\section{La couverture et ses grecques}

\begin{definition}[Couverture en delta d'une convertible]\label{def:s2:convertible-arbitrage:delta}
Une \emph{couverture en delta d'une convertible}\index{couverture en delta d'une convertible} est une position courte sur les actions de l'émetteur égale au delta de la convertible,
rééquilibrée à mesure que le delta varie, de sorte que la valeur de la position ne bouge pas avec de petites variations du cours de l'action et que son P\&L provient du gamma, du carry,
du crédit et de la décote de l'obligation.
\end{definition}

Agarwal, Fung, Loon et Naik ont montré qu'une stratégie d'achat et de couverture, longue de convertibles dont le risque action est couvert, explique une grande partie
des rendements des fonds. Les extrêmes de marché et l'offre de nouvelles convertibles comptent aussi. Ils interprètent les rendements des fonds comme la rémunération du financement
des émetteurs et du transfert d'une partie du risque action vers le marché des actions.

Mitchell, Pedersen et Pulvino décrivent la structure du marché. Les fonds d'arbitrage de convertibles détenaient jusqu'à 75~\% du marché des convertibles. Les convertibles
sont souvent émises en dessous de leur valeur fondamentale et continuent de se négocier en dessous, parce qu'elles sont peu liquides. Le risque principal de l'arbitragiste est que l'obligation
voie sa décote s'accroître davantage et doive être vendue au mauvais moment.

\texttt{firm.convarb} achète une obligation sur chacune des trente actions du marché synthétique à l'émission, au jour 1\,000, avec l'action au prix de conversion.
Les obligations s'achètent 3~\% sous la valeur du modèle, un écart qui tombe à 1~\% en six mois. Chaque position est courte de son delta, rééquilibrée chaque jour, avec des frais d'emprunt de 0.5~\% par
an. Le portefeuille est à effet de levier de trois pour un~: le capital représente un tiers de la valeur des obligations, soit 114.19 par obligation à l'achat. Le P\&L de chaque jour est réparti en catégories
(le \cref{lst:s2:convertible-arbitrage:book})~: le mouvement de l'obligation à crédit inchangé et la couverture en actions (ensemble, le gamma), la variation de l'état de
crédit, la variation de la décote, les coupons et l'emprunt.

\section{Crédit et emprunt}

\begin{definition}[Convertible couverte en crédit]\label{def:s2:convertible-arbitrage:credit}
Une \emph{convertible couverte en crédit}\index{convertible couverte en crédit} est une convertible couverte en delta à laquelle l'arbitragiste ajoute une protection de crédit
(credit default swaps, livre~2, chapitre~23, ou une position courte sur les obligations simples de l'émetteur) dimensionnée sur la sensibilité de l'obligation au crédit de l'émetteur,
de sorte que la position conserve sa valeur d'option et son carry sans le risque de défaut.
\end{definition}

La protection de crédit a un coût de prime. Dans le marché synthétique, elle rembourse la perte de l'obligation liée à toute hausse du multiplicateur de taux de défaut. Chaque année, elle coûte
la valeur en risque de l'obligation face à un taux de défaut doublé, divisée par l'annuité de la protection. L'emprunt constitue le coût de l'autre jambe. Une position courte verse au prêteur des
frais, et lorsque les actions sont rares, les frais augmentent ou le prêt fait l'objet d'un rappel.

\section{2005 et 2008}

Le marché synthétique prévoit un épisode 500 jours après l'achat, débutant avec le premier krach. Pendant trois mois, la décote des obligations passe de 1~\% à 6~\%
sous la valeur du modèle, le taux de défaut double et les frais d'emprunt montent à 5~\%. Ensuite, la décote revient à 1~\% en six mois. En capital~:

\begin{center}
\small
\begin{tabular}{@{}lcc@{}}
\toprule
trente obligations, trois pour un & couverte en delta & couverte en delta et en crédit\\
\midrule
carry hors épisode et reprise (\%/an) & 9.8 & 6.0\\
épisode de trois mois (\%) & $-11.6$ & $-5.8$\\
\quad dont décote & $-12.6$ & $-12.6$\\
\quad crédit, net de toute couverture & $-6.9$ & 0.0\\
\quad gamma (couverture de l'obligation et de l'action) & 3.7 & 3.7\\
reprise sur six mois (\%) & 22.0 & 10.6\\
cinq ans (\%) & 51.9 & 30.1\\
\bottomrule
\end{tabular}
\end{center}

Pendant l'épisode, les actions ont chuté avec le krach. Les obligations ont perdu 17.0~\% du capital à crédit inchangé, et l'action vendue à découvert a gagné 20.7~\%~: la
convexité a rapporté 3.7 points. Le creusement de la décote a coûté 12.6 points et le doublement du taux de défaut 6.9. Le coupon annuel qui a baissé pendant la fenêtre a ajouté 5.3, et
l'emprunt a coûté 1.1. La couverture de crédit a remboursé ses 6.9 points pendant l'épisode. Elle a coûté 21.8~\% du capital sur cinq ans, et elle a restitué sa part de
la reprise. Le portefeuille est long de tout à la fois (la \cref{fig:s2:convertible-arbitrage:book})~: volatilité, crédit et liquidité. Les pires jours pour chacun
sont les mêmes jours.

\begin{omfigure}
\begin{tikzpicture}
\begin{axis}[width=11cm,height=5.8cm,xlabel={\small années après l'achat},ylabel={\small cumulé, \% du capital},tick label style={font=\footnotesize},
  xmin=0,xmax=5,ymin=-20,ymax=60,grid=major,grid style={gray!25},table/col sep=comma,
  legend style={font=\scriptsize,at={(0.5,-0.3)},anchor=north},legend columns=5]
\addplot[black,thick] table[x=year,y=total]{figdata/strategies-2/08-convertible-arbitrage/book.csv};
\addplot[omThm,thick] table[x=year,y=cheapness]{figdata/strategies-2/08-convertible-arbitrage/book.csv};
\addplot[omDef,thick] table[x=year,y=credit]{figdata/strategies-2/08-convertible-arbitrage/book.csv};
\addplot[gray,thick,dashed] table[x=year,y=gamma]{figdata/strategies-2/08-convertible-arbitrage/book.csv};
\addplot[gray!60,thick] table[x=year,y=carry]{figdata/strategies-2/08-convertible-arbitrage/book.csv};
\legend{total,décote,crédit,gamma,coupons moins emprunt}
\end{axis}
\end{tikzpicture}

\omcaption{Le carnet synthétique de convertibles couvertes en delta, trente obligations à effet de levier trois pour un~: P\&L cumulé en part du capital et de ses compartiments sur
cinq ans. L'épisode commence à l'année 2.0 avec la première chute simulée. Données~: \texttt{s2\_convarb.cumulative}.}\label{fig:s2:convertible-arbitrage:book}
\end{omfigure}

Les épisodes réels obéissaient à leurs propres mécanismes. En 2005, la pression venait des investisseurs~: des rachats de plus de 20~\% du capital en un trimestre, 28
fonds spécialisés cédant 35~\% de leurs convertibles d'ici la fin de l'année, et une décote maximale autour des échéances de préavis de rachat. En
2008, selon Mitchell et Pulvino, la pression venait des prêteurs. La défaillance imminente des prime brokers a brutalement réduit l'effet de levier que les hedge funds
pouvaient obtenir, et ce qui semblait être un financement de long terme est devenu un financement de court terme. Dans les deux années, les arbitragistes qui auraient dû acheter les obligations bon marché étaient
ceux qui les vendaient.

\section{Fiches de stratégie}

\begin{strategyfile}{Convertible couverte en delta}\label{strat:s2:convertible-arbitrage:delta}
\sfield{Qui paie, et pourquoi} Les émetteurs, qui vendent des convertibles en dessous de leur valeur fondamentale pour obtenir rapidement des capitaux à bon marché~; les arbitragistes les financent et transfèrent le
risque actions au marché boursier.
\sfield{Instruments et plateformes} Obligations convertibles (surtout des émissions 144A)~; les actions des émetteurs, empruntées.
\sfield{Signal} Cours de marché comparé à la valeur du modèle~; la décote à l'émission.
\sfield{Dimensionnement et exécution} Couverte en delta, rééquilibrée au gré des variations du delta~; effet de levier fourni par les prime brokers.
\sfield{Coûts} Fourchettes des obligations, frais d'emprunt, financement.
\sfield{Comment elle meurt} Dépréciation par des ventes forcées~; retrait de l'effet de levier~; rappel des titres empruntés.
\sfield{Horizon, capacité, infrastructure} De quelques mois à plusieurs années~; un pricer de convertibles, un prêt de titres, un prime broker.
\sfield{Backtest honnête} Valeurs de modèle avec des entrées de crédit réalistes~; épisodes comme ceux de 2005 et 2008~; financement susceptible de disparaître.
\sfield{Sources} Agarwal, Fung, Loon et Naik (2011)~; Mitchell, Pedersen et Pulvino (2007)~: $-7.2~\%$ pour les fonds de janvier à mai 2005~; ce
chapitre~: 9.8~\% par an sur le capital, $-11.6~\%$ dans l'épisode.
\end{strategyfile}

\begin{strategyfile}{Convertible couverte en crédit}\label{strat:s2:convertible-arbitrage:credit}
\sfield{Qui paie, et pourquoi} Comme ci-dessus, moins la prime de crédit versée.
\sfield{Instruments et plateformes} Convertibles, actions, et credit default swaps (CDS) ou obligations classiques de l'émetteur.
\sfield{Signal} Décote comparée au coût de la protection.
\sfield{Dimensionnement et exécution} Protection dimensionnée selon la sensibilité de l'obligation au crédit, redimensionnée au gré des variations de l'action.
\sfield{Coûts} Primes de protection~; le spread du CDS.
\sfield{Comment elle meurt} La dépréciation n'est pas couverte~: elle a coûté 12.6 points dans l'épisode du chapitre, dans les deux cas.
\sfield{Horizon, capacité, infrastructure} Comme ci-dessus~; lignes de trading de crédit.
\sfield{Backtest honnête} Imputer chaque jour le coût de la protection, et non seulement quand elle rapporte.
\sfield{Sources} Ce chapitre~: carry de 6.0~\% au lieu de 9.8~\%, l'épisode à $-5.8~\%$ au lieu de $-11.6~\%$.
\end{strategyfile}

\begin{strategyfile}{Opération de crédit sur convertible «~busted~»}\label{strat:s2:convertible-arbitrage:busted}
\sfield{Qui paie, et pourquoi} Les détenteurs contraints de vendre des convertibles qui ne se comportent plus comme des actions, et les acheteurs de crédit qui ne peuvent pas les détenir.
\sfield{Instruments et plateformes} \omterm{def:s2:convertible-arbitrage:busted}{Convertibles «~busted~»}~; la dette classique de l'émetteur ou ses CDS.
\sfield{Signal} Le rendement de la convertible comparé au crédit classique de l'émetteur.
\sfield{Dimensionnement et exécution} Dimensionnée au crédit~; petites couvertures en actions issues du delta du modèle action-crédit.
\sfield{Coûts} Obligations peu liquides.
\sfield{Comment elle meurt} Défaut~: l'obligation verse son taux de recouvrement.
\sfield{Horizon, capacité, infrastructure} Quelques mois~; analyse de crédit.
\sfield{Backtest honnête} Recouvrements et restructurations tels qu'ils se sont produits.
\sfield{Sources} Aucun chiffre de performance vérifié.
\end{strategyfile}

\begin{strategyfile}{Convertible à l'émission}\label{strat:s2:convertible-arbitrage:newissue}
\sfield{Qui paie, et pourquoi} Les émetteurs qui acceptent une décote pour aller vite~: les convertibles sont souvent émises en dessous de leur valeur fondamentale.
\sfield{Instruments et plateformes} Nouvelles émissions 144A~; les actions des émetteurs.
\sfield{Signal} Prix d'émission comparé à la valeur du modèle.
\sfield{Dimensionnement et exécution} Allocations lors du book-build~; couvertes dès le premier jour.
\sfield{Coûts} Emprunt sur les actions nouvellement vendues à découvert.
\sfield{Comment elle meurt} Nouvelles émissions surpeuplées~; décotes qui se réduisent.
\sfield{Horizon, capacité, infrastructure} Quelques semaines~; relations avec les souscripteurs.
\sfield{Backtest honnête} Les allocations n'étaient pas garanties~; supposer des exécutions partielles.
\sfield{Sources} Mitchell, Pedersen et Pulvino (2007)~; la décote de 3~\% à l'émission de ce chapitre est une hypothèse.
\end{strategyfile}

\section{Tutoriel~: long sur tout à la fois}\label{tut:s2:convertible-arbitrage:tut}

\begin{tutorial}
\textbf{Objectif.} Pricer des convertibles avec le pricer du livre~5, les couvrir en delta et en crédit, et faire passer le carnet dans l'épisode prévu.
\textbf{État final~:} le tableau et les deux figures.
\begin{tutsteps}
\item \textbf{Le carnet}.
\omcode{code/firm/convarb/firm_convarb.py}{104}{135}{Une obligation, vendue à découvert sur son delta, avec protection de crédit facultative, attribuée jour par jour.}\label{lst:s2:convertible-arbitrage:book}
\item \textbf{Rendements sur capital}.
\omcode{code/strategies-2/08-convertible-arbitrage/python/s2_convarb.py}{46}{60}{Carry, épisode par compartiment, recouvrement.}\label{lst:s2:convertible-arbitrage:summary}
\item \textbf{Exécuter} \texttt{summary()}, \texttt{summary(True)}, \texttt{profile()} et \texttt{fig\_convarb.py}.
\end{tutsteps}
\textbf{Pour aller plus loin.} Faire en sorte que le rappel de titres empruntés force un débouclage partiel de la couverture~; porter l'effet de levier à cinq pour un et trouver la perte d'épisode qui
déclenche un appel de marge~; démarrer le carnet à un cours d'action de 20 et voir de quoi est faite la couverture d'une \omterm{def:s2:convertible-arbitrage:busted}{convertible «~busted~»}.
\end{tutorial}

\section{Construction~: arbitrage de convertibles}\label{bld:s2:convertible-arbitrage:convarb}

\begin{build}
\bfield{Rôle} Un carnet de convertibles sur le pricer du livre~5, avec couvertures en delta et en crédit, emprunt, et un scénario de ventes forcées.
\bfield{Interface} \texttt{ArbConfig(\dots)}, \texttt{value\_tables(cfg)}, \texttt{mark(tables, tau, S, stressed)},
\texttt{delta(tables, tau, S, stressed)}, \texttt{scenario(T, cfg)}, \texttt{run\_book(S, cfg, credit\_hedge)}.
\bfield{Règles} Valeurs tabulées mensuellement par maturité pour une intensité de défaut normale et une intensité stressée~; le cours de marché égale la valeur du modèle moins la décote~; les compartiments
s'additionnent pour donner le total.
\bfield{Tests d'acceptation} \texttt{code/firm/convarb/tests/}~: la forme du scénario~; les valorisations égalent le pricer aux nœuds de la table, le stress fait baisser la valeur,
le delta est inférieur au ratio de conversion~; les compartiments s'additionnent sur un chemin calme, sans P\&L de crédit.
\bfield{Pour aller plus loin} Rappels de titres empruntés~; appels de marge~; corrélation du crédit de plusieurs émetteurs.
\end{build}

\begin{omsources}
\item M.~Mitchell, L.~H.~Pedersen et T.~Pulvino, ``Slow moving capital'', \emph{American Economic Review} 97(2), 2007 (document de travail NBER 12877).
\item M.~Mitchell et T.~Pulvino, ``Arbitrage crashes and the speed of capital'', \emph{Journal of Financial Economics} 104(3), 2012.
\item V.~Agarwal, W.~H.~Fung, Y.~C.~Loon et N.~Y.~Naik, ``Risk and return in convertible arbitrage: evidence from the convertible bond market'',
\emph{Journal of Empirical Finance} 18(2), 2011.
\end{omsources}

\section{Exercices}

\begin{exercise}[$\star$]\label{exo:s2:convertible-arbitrage:1}
Une obligation se convertit en 2.5 actions. Calculer son prix de conversion pour 100 de nominal, ainsi que sa valeur de conversion à un cours d'action de 60.
\end{exercise}

\begin{exercise}[$\star$]\label{exo:s2:convertible-arbitrage:2}
Les obligations se déprécient de 2.7~\%. Calculer la perte sur son capital d'un fonds à effet de levier trois pour un.
\end{exercise}

\begin{exercise}[$\star$]\label{exo:s2:convertible-arbitrage:3}
Une obligation valant 81.2 a un delta de 1.27 action pour un cours d'action de 15. Calculer son élasticité aux actions.
\end{exercise}

\begin{exercise}[$\star\star$]\label{exo:s2:convertible-arbitrage:4}
Expliquer pourquoi une \omterm{def:s2:convertible-arbitrage:busted}{convertible «~busted~»} peut valoir moins que l'obligation nue au spread de crédit de l'émission.
\end{exercise}

\begin{exercise}[$\star\star$]\label{exo:s2:convertible-arbitrage:5}
Dans l'épisode, les obligations ont perdu 17.0~\% de capital et la couverture par l'action a gagné 20.7~\%. Expliquer pourquoi la couverture a gagné davantage que les obligations n'ont perdu.
\end{exercise}

\begin{exercise}[$\star\star$]\label{exo:s2:convertible-arbitrage:6}
Expliquer pourquoi le risque d'élargissement du spread est un risque qu'aucune couverture ne supprime.
\end{exercise}

\begin{exercise}[$\star\star\star$]\label{exo:s2:convertible-arbitrage:7}
\emph{Programmation.} Comparer \texttt{summary()} et \texttt{summary(True)}. Déterminer combien coûte la couverture en crédit sur cinq ans, et ce qu'elle rapporte.
\end{exercise}

\begin{exercise}[$\star\star\star$]\label{exo:s2:convertible-arbitrage:8}
\emph{Trouver la faille.} «~Notre livre de convertibles est couvert en delta et en crédit, il n'a donc aucun risque et nous pouvons le lever à dix contre un.~»
\end{exercise}

\section{Problème~: tout acheter en même temps}

\begin{problem}[{Problème du week-end --- un livre d'arbitrage de convertibles}]\label{pb:s2:convertible-arbitrage:1}
Le livre synthétique du chapitre, le pricer du livre~5 et les données publiques.

\medskip\noindent\textbf{Partie I --- L'obligation.}
\begin{enumerate}
\item Que détient le porteur d'une convertible~?
\item Donner la valeur de l'obligation du chapitre, sa valeur de conversion et son delta à l'émission.
\item Définir une \omterm{def:s2:convertible-arbitrage:busted}{convertible «~busted~»} et décrire l'obligation pour un cours d'action de 15.
\item Pourquoi le lien action-crédit importe-t-il pour la couverture~?
\end{enumerate}

\medskip\noindent\textbf{Partie II --- La transaction.}
\begin{enumerate}[resume]
\item Définir la \omterm{def:s2:convertible-arbitrage:delta}{couverture en delta d'une convertible} et la \omterm{def:s2:convertible-arbitrage:credit}{convertible couverte en crédit}.
\item Que trouvent Agarwal et ses coauteurs~?
\item Que trouvent Mitchell, Pedersen et Pulvino au sujet de 2005~?
\item Qu'est-ce qui change en 2008~?
\end{enumerate}

\medskip\noindent\textbf{Partie III --- Le livre.}
\begin{enumerate}[resume]
\item Comment le livre synthétique est-il construit, et comment son P\&L est-il attribué~?
\item Donner son carry, avec et sans la couverture en crédit.
\item Décomposer la perte de l'épisode.
\item Que se passe-t-il pendant la reprise~?
\end{enumerate}

\medskip\noindent\textbf{Partie IV --- Le verdict.}
\begin{enumerate}[resume]
\item Énoncer le \emph{résultat nommé}~: le carry du livre couvert et sa perte dans l'épisode de vente forcée programmée, ainsi que la part due au crédit.
\item Que coûte et que rapporte la couverture en crédit~?
\item Pourquoi l'effet de levier est-il le vrai risque du livre~?
\item Pourquoi les arbitragistes étaient-ils vendeurs quand les obligations étaient les moins chères~?
\item Comment dimensionner le livre~?
\item Comment mener un backtest honnête de celui-ci~?
\item Quelle fiche de stratégie dépend le plus de la décote à l'émission~?
\item En une phrase~: que rémunère un arbitragiste de convertibles~?
\end{enumerate}
\end{problem}

\section{Questions d'entretien}

\begin{interviewq}[$\star$ \iqroles{trader}]\label{iq:s2:convertible-arbitrage:1}
Décrivez étape par étape une transaction d'arbitrage de convertibles.
\end{interviewq}

\begin{interviewq}[$\star\star$ \iqroles{chercheur}]\label{iq:s2:convertible-arbitrage:2}
Comment le lien entre le cours de l'action et le crédit modifie-t-il le delta d'une convertible~?
\end{interviewq}

\begin{interviewq}[$\star\star$ \iqroles{risque}]\label{iq:s2:convertible-arbitrage:3}
Un livre de convertibles est couvert en delta et en crédit. Quels risques subsistent, et comment les soumettriez-vous à des tests de résistance~?
\end{interviewq}

\begin{interviewq}[$\star\star$ \iqroles{trader}]\label{iq:s2:convertible-arbitrage:4}
L'emprunt de titres sur l'une de vos couvertures courtes est rappelé. Que faites-vous~?
\end{interviewq}

\begin{interviewq}[$\star\star$ \iqroles{développeur}]\label{iq:s2:convertible-arbitrage:5}
Concevez le processus quotidien qui valorise et couvre un livre de 200 convertibles.
\end{interviewq}

\begin{interviewq}[$\star\star\star$ \iqroles{chercheur}]\label{iq:s2:convertible-arbitrage:6}
Montrez qu'une convertible couverte en delta rapporte environ $\tfrac12\Gamma S^2(\sigma_r^2 - \sigma_i^2)\,\mathrm{d}t$ plus le carry, et expliquez ce que fait un
défaut à la couverture.
\end{interviewq}
